\documentclass{article}

\usepackage{bm}
\usepackage{ctex}
\usepackage{tikz}
\usepackage{float}
\usepackage{xcolor}
\usepackage{dsfont}
\usepackage{setspace}
\usepackage{datetime}
\usepackage{geometry}
\usepackage{graphicx}
\usepackage{enumitem}
\usepackage{tcolorbox}
\usepackage{nicematrix}
\usepackage{algorithm, algpseudocode}
\usepackage{amsmath, amssymb, amsfonts, mathrsfs}
\usepackage[fontsize = 12pt]{fontsize}

\geometry{
    a4paper,
    left = 2.1cm,
    right = 2.1cm,
    top = 2.5cm,
    bottom = 2.5cm
}

\setitemize[1]{
    itemsep = 7pt,
    partopsep = 0pt,
    parsep = \parskip,
    topsep = 5pt
}

\renewcommand{\thesubsection}{\alph{subsection})}
\newcommand{\diff}[2]{\dfrac{\mathrm{d}#1}{\mathrm{d}#2}}
\newcommand{\diffn}[3]{\dfrac{\mathrm{d}^{#3}#1}{\mathrm{d}#2^{#3}}}
\newcommand{\piff}[2]{\dfrac{\partial{#1}}{\partial{#2}}}
\newcommand{\eval}[2]{\left.#1\right|_{#2}}
\newcommand{\e}{\mathrm{e}}
\renewcommand{\d}{\mathrm{d}}
\newcommand{\barline}[1]{\overline{#1\rule{0pt}{1.75ex}}}

\newtheorem{theorem}{定理}

\title{格密码：第二次作业}
\author{Illusionna \quad 2025..............004 \quad 人工智能学院}
\date{\currenttime,\ \today}



\begin{document}

\maketitle
\thispagestyle{empty}
\begin{spacing}{2}
    \tableofcontents
\end{spacing}
\thispagestyle{empty}
\clearpage
\setcounter{page}{1}



\section{给出下面基的 3/4-LLL 约化基}
\[ \setlength{\arraycolsep}{7pt}\begin{pmatrix}
    1 & -1 & 3 \\
    1 & 0 & 5 \\
    1 & 2 & 6 \\
\end{pmatrix} \]

LLL 算法中，一组基是 LLL 约化的，必须同时满足两个条件：

\begin{itemize}
    \item 长度约化条件：对于所有 $j<i$，正交化系数必须满足 $|\mu_{i,j}|\leqslant0.5$
    \item Lov\'asz 交换条件：相邻基向量之间满足特定长度递减界限 $\|\tilde{\bm{b}}_{k}\|^2\geqslant(\delta-\mu_{k,k-1}^2)\|\tilde{\bm{b}}_{k-1}\|^2$
\end{itemize}

参数 $\delta=3/4$，且基向量：
\[ \bm{b}_1=(1,-1,3)\qquad\bm{b}_2=(1,0,5)\qquad\bm{b}_3=(1,2,6) \]

\paragraph*{第 1 轮}~{}

计算 Gram-Schmidt 正交化：
\[ \tilde{\bm{b}}_1=\bm{b}_1=(1,-1,3)\qquad \|\tilde{\bm{b}}_1\|^2=11 \]

计算投影系数：
\[ \mu_{2,1}=\dfrac{\bm{b}_2\cdot\tilde{\bm{b}}_1}{\|\tilde{\bm{b}}_1\|^2}=\dfrac{1+0+15}{11}\approx1.4545 \]

由于 $|\mu_{2,1}|>\dfrac{1}{2}$，所以对 $\bm{b}_2$ 进行长度约化：
\[ \bm{b}_2\leftarrow\bm{b}_2-\left\lfloor\mu_{2,1}\right\rceil\cdot\bm{b}_1=(0,1,2) \]

更新投影系数：
\[ \mu_{2,1}=\dfrac{\bm{b}_2\cdot\tilde{\bm{b}}_1}{\|\tilde{\bm{b}}_1\|^2}=\dfrac{5}{11} \]

则：
\[ \tilde{\bm{b}}_2=\bm{b}_2-\mu_{2,1}\tilde{\bm{b}}_1=\left(-\dfrac{5}{11},\ \dfrac{16}{11},\ \dfrac{7}{11}\right)\qquad\|\tilde{\bm{b}}_2\|^2=\dfrac{30}{11}\approx2.7273 \]

检查到 Lov\'asz 条件不成立：
\[ \|\tilde{\bm{b}}_{2}\|^2\approx2.7273>(\delta-\mu_{2,1}^2)\|\tilde{\bm{b}}_{1}\|^2=\left(\dfrac{3}{4}-\dfrac{25}{121}\right)\times11\approx5.9773 \]

所以交换：
\[ \bm{b}_1\leftrightarrow\bm{b}_2 \]

因此：
\[ \bm{b}_1=(0,1,2)\qquad\bm{b}_2=(1,-1,3)\qquad\bm{b}_3=(1,2,6) \]



\paragraph*{第 2 轮}~{}

计算 Gram-Schmidt 正交化：
\[ \tilde{\bm{b}}_1=\bm{b}_1=(0,1,2)\qquad \|\tilde{\bm{b}}_1\|^2=5 \]

计算投影系数：
\[ \mu_{2,1}=\dfrac{\bm{b}_2\cdot\tilde{\bm{b}}_1}{\|\tilde{\bm{b}}_1\|^2}=\dfrac{0-1+6}{5}=1 \]

由于 $|\mu_{2,1}|>\dfrac{1}{2}$，所以对 $\bm{b}_2$ 进行长度约化：
\[ \bm{b}_2\leftarrow\bm{b}_2-\left\lfloor\mu_{2,1}\right\rceil\cdot\bm{b}_1=(1,-2,1) \]

更新投影系数：
\[ \mu_{2,1}=\dfrac{\bm{b}_2\cdot\tilde{\bm{b}}_1}{\|\tilde{\bm{b}}_1\|^2}=0 \]

则：
\[ \tilde{\bm{b}}_2=\bm{b}_2-\mu_{2,1}\tilde{\bm{b}}_1=(1,-2,1)\qquad\|\tilde{\bm{b}}_2\|^2=6 \]

检查到 Lov\'asz 条件成立：
\[ \|\tilde{\bm{b}}_{2}\|^2=6>(\delta-\mu_{2,1}^2)\|\tilde{\bm{b}}_{1}\|^2=\left(\dfrac{3}{4}-0\right)\times5=3.75 \]

因此：
\[ \bm{b}_1=(0,1,2)\qquad\bm{b}_2=(1,-2,1)\qquad\bm{b}_3=(1,2,6) \]



\paragraph*{第 3 轮}~{}

计算 Gram-Schmidt 正交化：
\[ \tilde{\bm{b}}_1=\bm{b}_1=(0,1,2)\qquad \|\tilde{\bm{b}}_1\|^2=5 \]
\[ \tilde{\bm{b}}_2=\bm{b}_2-\mu_{2,1}\tilde{\bm{b}}_1=(1,-2,1)\qquad\|\tilde{\bm{b}}_2\|^2=6 \]

引入 $\bm{b}_3=(1,2,6)$，则有：
\[ \mu_{3,1}=\dfrac{\bm{b}_3\cdot\tilde{\bm{b}}_1}{\|\tilde{\bm{b}}_1\|^2}=\dfrac{14}{5}=2.8\qquad\mu_{3,2}=\dfrac{\bm{b}_3\cdot\tilde{\bm{b}}_2}{\|\tilde{\bm{b}}_2\|^2}=\dfrac{3}{6}=0.5 \]

\begin{itemize}
    \item $j=2$ 时，$|\mu_{3,2}|=\dfrac{1}{2}$ 满足条件，所以不做强制约化
    \[ \bm{b}_3\leftarrow\bm{b}_3-\left\lfloor\mu_{3,2}\right\rceil\bm{b}_2=\bm{b}_3-\begin{cases}
        0\cdot\bm{b}_2=(1,2,6)
        \\[7pt]
        1\cdot\bm{b}_2=(0,4,5)
    \end{cases}\qquad\mu_{3,2}=\begin{cases}
        +0.5
        \\[7pt]
        -0.5
    \end{cases} \]
    \item $j=1$ 时，$|\mu_{3,1}|>\dfrac{1}{2}$ 不满足，所以对 $\bm{b}_3$ 进行长度约化：
    \[ \bm{b}_3\leftarrow\bm{b}_3-\left\lfloor\mu_{3,1}\right\rceil\bm{b}_1=(1,-1,0)\qquad\mu_{3,1}=\dfrac{\bm{b}_3\cdot\tilde{\bm{b}}_1}{\|\tilde{\bm{b}}_1\|^2}=-0.2 \]
\end{itemize}

则：
\[ \tilde{\bm{b}}_3=\bm{b}_3-\mu_{3,1}\tilde{\bm{b}}_1-\mu_{3,2}\tilde{\bm{b}}_2=(0.5,\ 0.2,\ -0.1)\qquad\|\tilde{\bm{b}}_3\|^2=0.3 \]

检查到 Lov\'asz 条件不成立：
\[ \|\tilde{\bm{b}}_{3}\|^2=0.3>(\delta-\mu_{3,2}^2)\|\tilde{\bm{b}}_{2}\|^2=\left(\dfrac{3}{4}-0.25\right)\times6=3 \]

所以交换：
\[ \bm{b}_2\leftrightarrow\bm{b}_3 \]

因此：
\[ \bm{b}_1=(0,1,2)\qquad\bm{b}_2=(1,-1,0)\qquad\bm{b}_3=(1,-2,1) \]



\paragraph*{第 4 轮}~{}

计算 Gram-Schmidt 正交化：
\[ \tilde{\bm{b}}_1=\bm{b}_1=(0,1,2)\qquad \|\tilde{\bm{b}}_1\|^2=5 \]

计算投影系数：
\[ \mu_{2,1}=\dfrac{\bm{b}_2\cdot\tilde{\bm{b}}_1}{\|\tilde{\bm{b}}_1\|^2}=\dfrac{0-1+0}{5}=-0.2 \]

由于 $|\mu_{2,1}|\leqslant\dfrac{1}{2}$，无需减去整数倍的 $\bm{b}_1$，计算 Gram-Schmidt 正交化：
\[ \tilde{\bm{b}}_2=\bm{b}_2-\mu_{2,1}\tilde{\bm{b}}_1=(1,-0.8,0.4)\qquad\|\tilde{\bm{b}}_2\|^2=1.8 \]

检查到 Lov\'asz 条件不成立：
\[ \|\tilde{\bm{b}}_{2}\|^2=1.8>(\delta-\mu_{2,1}^2)\|\tilde{\bm{b}}_{1}\|^2=\left(\dfrac{3}{4}-0.04\right)\times5=3.55 \]

所以交换：
\[ \bm{b}_1\leftrightarrow\bm{b}_2 \]

因此：
\[ \bm{b}_1=(1,-1,0)\qquad\bm{b}_2=(0,1,2)\qquad\bm{b}_3=(1,-2,1) \]



\paragraph*{第 5 轮}~{}

计算 Gram-Schmidt 正交化：
\[ \tilde{\bm{b}}_1=\bm{b}_1=(1,-1,0)\qquad \|\tilde{\bm{b}}_1\|^2=2 \]

计算投影系数：
\[ \mu_{2,1}=\dfrac{\bm{b}_2\cdot\tilde{\bm{b}}_1}{\|\tilde{\bm{b}}_1\|^2}=\dfrac{0-1+0}{2}=-0.5 \]

由于 $|\mu_{2,1}|=0.5$，所以不进行强制约化：
\[ \bm{b}_2\leftarrow\bm{b}_2-\left\lfloor\mu_{2,1}\right\rceil\bm{b}_1=\bm{b}_2-\begin{cases}
    0\cdot\bm{b}_1=(0,1,2)
    \\[7pt]
    -1\cdot\bm{b}_1=(1,0,2)
\end{cases}\qquad\mu_{2,1}=\begin{cases}
    -0.5
    \\[7pt]
    +0.5
\end{cases} \]

计算 Gram-Schmidt 正交化：
\[ \tilde{\bm{b}}_2=\bm{b}_2-\mu_{2,1}\tilde{\bm{b}}_1=\begin{cases}
    (0,1,2)-(-0.5)\cdot(1,-1,0)=(0.5,0.5,2)\qquad\|\tilde{\bm{b}}_2\|^2=4.5
    \\[7pt]
    (1,0,2)-0.5\cdot(1,-1,0)=(0.5,0.5,2)\qquad\|\tilde{\bm{b}}_2\|^2=4.5
\end{cases} \]

因此，保留 $\bm{b}_2$ 不变，检查到 Lov\'asz 条件成立：
\[ \|\tilde{\bm{b}}_{2}\|^2=4.5>(\delta-\mu_{2,1}^2)\|\tilde{\bm{b}}_{1}\|^2=\left(\dfrac{3}{4}-0.25\right)\times2=1 \]

所以：
\[ \bm{b}_1=(1,-1,0)\qquad\bm{b}_2=(0,1,2)\qquad\bm{b}_3=(1,-2,1) \]



\paragraph*{第 6 轮}~{}

计算 Gram-Schmidt 正交化：
\[ \tilde{\bm{b}}_1=\bm{b}_1=(1,-1,0)\qquad \|\tilde{\bm{b}}_1\|^2=2 \]
\[ \tilde{\bm{b}}_2=\bm{b}_2-\mu_{2,1}\tilde{\bm{b}}_1=(0.5,0.5,2)\qquad\|\tilde{\bm{b}}_2\|^2=4.5 \]

引入 $\bm{b}_3=(1,-2,1)$，则有：
\[ \mu_{3,1}=\dfrac{\bm{b}_3\cdot\tilde{\bm{b}}_1}{\|\tilde{\bm{b}}_1\|^2}=\dfrac{1+2+0}{2}=1.5\qquad\mu_{3,2}=\dfrac{\bm{b}_3\cdot\tilde{\bm{b}}_2}{\|\tilde{\bm{b}}_2\|^2}=\dfrac{0.5-1+2}{4.5}\approx0.3333 \]

\begin{itemize}
    \item 对于 $j=2$ 时，$|\mu_{3,2}|\leqslant\dfrac{1}{2}$，不做强制约化，跳过
    \item 对于 $j=1$ 时，$|\mu_{3,1}|>\dfrac{1}{2}$，进行约化：
    \[ \bm{b}_3\leftarrow\bm{b}_3-\left\lfloor\mu_{3,1}\right\rceil\bm{b}_1=\bm{b}_3-\begin{cases}
        2\cdot\bm{b}_1=(-1,0,1)\qquad\mu_{3,1}=\dfrac{\bm{b}_3\cdot\tilde{\bm{b}}_1}{\|\tilde{\bm{b}}_1\|^2}=-0.5
        \\[10pt]
        1\cdot\bm{b}_1=(0,-1,1)\qquad\mu_{3,1}=\dfrac{\bm{b}_3\cdot\tilde{\bm{b}}_1}{\|\tilde{\bm{b}}_1\|^2}=+0.5
    \end{cases} \]
\end{itemize}

保留 $\mu_{3,1}=-0.5$，计算 Gram-Schmidt 正交化：
\[ \tilde{\bm{b}}_3=\bm{b}_3-\mu_{3,1}\tilde{\bm{b}}_1-\mu_{3,2}\tilde{\bm{b}}_2=\left(-\dfrac{2}{3},\ -\dfrac{2}{3},\ \dfrac{1}{3}\right)\qquad\|\tilde{\bm{b}}_3\|^2=1 \]

检查到 Lov\'asz 条件不成立：
\[ \|\tilde{\bm{b}}_{3}\|^2=1>(\delta-\mu_{3,2}^2)\|\tilde{\bm{b}}_{2}\|^2=\left(\dfrac{3}{4}-\dfrac{1}{9}\right)\times4.5=2.875 \]

所以交换：
\[ \bm{b}_2\leftrightarrow\bm{b}_3 \]

因此：
\[ \bm{b}_1=(1,-1,0)\qquad\bm{b}_2=(-1,0,1)\qquad\bm{b}_3=(0,1,2) \]

\paragraph*{第 7 轮}~{}

计算 Gram-Schmidt 正交化：
\[ \tilde{\bm{b}}_1=\bm{b}_1=(1,-1,0)\qquad \|\tilde{\bm{b}}_1\|^2=2 \]

计算投影系数：
\[ \mu_{2,1}=\dfrac{\bm{b}_2\cdot\tilde{\bm{b}}_1}{\|\tilde{\bm{b}}_1\|^2}=\dfrac{-1+0+0}{2}=-0.5 \]

由于 $|\mu_{2,1}|\leqslant\dfrac{1}{2}$，不做约化，计算 Gram-Schmidt 正交化：
\[ \tilde{\bm{b}}_2=\bm{b}_2-\mu_{2,1}\tilde{\bm{b}}_1=(-0.5,-0.5,1)\qquad\|\tilde{\bm{b}}_2\|^2=1.5 \]

检查到 Lov\'asz 条件成立：
\[ \|\tilde{\bm{b}}_{2}\|^2=1.5>(\delta-\mu_{2,1}^2)\|\tilde{\bm{b}}_{1}\|^2=\left(\dfrac{3}{4}-0.25\right)\times2=1 \]

因此：
\[ \bm{b}_1=(1,-1,0)\qquad\bm{b}_2=(-1,0,1)\qquad\bm{b}_3=(0,1,2) \]



\paragraph*{第 8 轮}~{}

计算 Gram-Schmidt 正交化：
\[ \tilde{\bm{b}}_1=\bm{b}_1=(1,-1,0)\qquad \|\tilde{\bm{b}}_1\|^2=2 \]
\[ \tilde{\bm{b}}_2=\bm{b}_2-\mu_{2,1}\tilde{\bm{b}}_1=(-0.5,-0.5,1)\qquad\|\tilde{\bm{b}}_2\|^2=1.5 \]

引入 $\bm{b}_3=(0,1,2)$，则有：
\[ \mu_{3,1}=\dfrac{\bm{b}_3\cdot\tilde{\bm{b}}_1}{\|\tilde{\bm{b}}_1\|^2}=\dfrac{0-1+0}{2}=-0.5\qquad\mu_{3,2}=\dfrac{\bm{b}_3\cdot\tilde{\bm{b}}_2}{\|\tilde{\bm{b}}_2\|^2}=\dfrac{0-0.5+2}{1.5}=1 \]

由于 $|\mu_{3,2}|>\dfrac{1}{2}$，所以对 $\bm{b}_3$ 进行长度约化：
\[ \bm{b}_3\leftarrow\bm{b}_3-\left\lfloor\mu_{3,2}\right\rceil\bm{b}_2=(1,1,1) \]

更新投影系数：
\[ \mu_{3,1}=\dfrac{\bm{b}_3\cdot\tilde{\bm{b}}_1}{\|\tilde{\bm{b}}_1\|^2}=0\qquad\mu_{3,2}=\dfrac{\bm{b}_3\cdot\tilde{\bm{b}}_2}{\|\tilde{\bm{b}}_2\|^2}=0 \]

约化非常彻底，计算 Gram-Schmidt 正交化：
\[ \tilde{\bm{b}}_3=\bm{b}_3-\mu_{3,1}\tilde{\bm{b}}_1-\mu_{3,2}\tilde{\bm{b}}_2=(1,1,1)\qquad\|\tilde{\bm{b}}_3\|^2=3 \]

检查到 Lov\'asz 条件成立：
\[ \|\tilde{\bm{b}}_{3}\|^2=3>(\delta-\mu_{3,2}^2)\|\tilde{\bm{b}}_{2}\|^2=\left(\dfrac{3}{4}-0\right)\times1.5=1.125 \]

因此：
\[ \bm{b}_1=(1,-1,0)\qquad\bm{b}_2=(-1,0,1)\qquad\bm{b}_3=(1,1,1) \]



\paragraph*{迭代完成，最终收敛，超出了基向量的数量}~{}
\[ B_{\text{LLL}}=\begin{pmatrix}
    1 & -1 & 0 \\
    -1 & 0 & 1 \\
    1 & 1 & 1 \\
\end{pmatrix} \]

验证 $|\mu_{i,j}|\leqslant\dfrac{1}{2}$ 均成立：
\[ \mu_{2,1}=\dfrac{\bm{b}_2\cdot\tilde{\bm{b}}_1}{\|\tilde{\bm{b}}_1\|^2}=-0.5\qquad\mu_{3,1}=\dfrac{\bm{b}_3\cdot\tilde{\bm{b}}_1}{\|\tilde{\bm{b}}_1\|^2}=0\qquad\mu_{3,2}=\dfrac{\bm{b}_3\cdot\tilde{\bm{b}}_2}{\|\tilde{\bm{b}}_2\|^2}=0 \]

验证 Lov\'asz 条件 $\|\tilde{\bm{b}}_{k}\|^2\geqslant(\delta-\mu_{k,k-1}^2)\|\tilde{\bm{b}}_{k-1}\|^2$ 均成立：
\begin{align*}
    \|\tilde{\bm{b}}_{3}\|^2&=3>(\delta-\mu_{3,2}^2)\|\tilde{\bm{b}}_{2}\|^2=\left(\dfrac{3}{4}-0\right)\times1.5=1.125
    \\
    \|\tilde{\bm{b}}_{2}\|^2&=1.5>(\delta-\mu_{2,1}^2)\|\tilde{\bm{b}}_{1}\|^2=\left(\dfrac{3}{4}-0.25\right)\times2=1
\end{align*}



\section{证明 3/4-LLL 约化基的下面性质}
\subsection{$\|\bm{b}_1\|\leqslant2^{(n-1)/4}(\det\Lambda)^{1/n}$}

\begin{theorem}
    求和
    \[ \sum_{i=1}^{n}(i-1)=\dfrac{n(n-1)}{2} \]
\end{theorem}

\begin{theorem}
    求积
    \[ \prod_{i=1}^{n}\|\tilde{\bm{b}}_{i}\|^2=(\det\Lambda)^2 \]
\end{theorem}

因为：
\[ \|\bm{b}_1\|^2\leqslant2^{i-1}\|\tilde{\bm{b}}_{i}\|^2 \]

所以：
\begin{align*}
    \prod_{i=1}^{n}\|\bm{b}_1\|^2&\leqslant\prod_{i=1}^{n}2^{i-1}\|\tilde{\bm{b}}_{i}\|^2
    \\
    \|\bm{b}_1\|^{2n}&\leqslant2^{n(n-1)/2}(\det\Lambda)^2
\end{align*}

两边同时开 $2n$ 次方，即：
\[ \|\bm{b}_1\|\leqslant2^{(n-1)/4}(\det\Lambda)^{1/n} \]



\subsection{对任意 $1\leqslant i\leqslant n,\ \|\bm{b}_i\|\leqslant2^{(i-1)/2}\|\tilde{\bm{b}}_i\|$}

\begin{theorem}
    Gram-Schmidt 正交化
    \[ \bm{b}_i=\tilde{\bm{b}}_i+\sum_{j=1}^{i-1}\mu_{i,j}\tilde{\bm{b}}_j\qquad\|\mu_{i,j}\|\leqslant\dfrac{1}{2} \]
\end{theorem}

因为：
\[ \|\bm{b}_i\|^2=\|\tilde{\bm{b}}_i\|^2+\sum_{j=1}^{i-1}\mu_{i,j}^2\|\tilde{\bm{b}}_j\|^2\leqslant\|\tilde{\bm{b}}_i\|^2+\dfrac{1}{4}\sum_{j=1}^{i-1}2^{i-j}\|\tilde{\bm{b}}_i\|^2 \]

所以：
\[ \|\bm{b}_i\|^2\leqslant(1+2^{i-2}-0.5)\|\bm{b}_i\|^2\leqslant2^{i-1}\|\bm{b}_i\|^2 \]

开方即可得：
\[ \|\bm{b}_i\|\leqslant2^{(i-1)/2}\|\tilde{\bm{b}}_i\| \]



\subsection{$\displaystyle\prod\|\bm{b}_i\|\leqslant2^{n(n-1)/4}\det\Lambda$}
因为对任意 $1\leqslant i\leqslant n$ 有：
\[ \|\bm{b}_i\|\leqslant2^{(i-1)/2}\|\tilde{\bm{b}}_i\| \]

所以：
\[ \prod_{i=1}^{n}\|\bm{b}_i\|\leqslant\prod_{i=1}^{n}\left[2^{(i-1)/2}\|\tilde{\bm{b}}_i\|\right]=2^{\sum\limits_{i=1}^{n}\tfrac{i-1}{2}}\cdot\prod_{i=1}^{n}\|\tilde{\bm{b}}_i\|=2^{n(n-1)/4}\det\Lambda \]



\subsection{对任意 $1\leqslant i\leqslant j\leqslant n,\ \|\bm{b}_i\|\leqslant2^{(j-1)/2}\|\tilde{\bm{b}}_j\|$}
因为对任意 $1\leqslant i\leqslant n$ 有：
\[ \|\bm{b}_i\|\leqslant2^{(i-1)/2}\|\tilde{\bm{b}}_i\|\implies\|\bm{b}_i\|^2\leqslant2^{i-1}\|\tilde{\bm{b}}_i\|^2 \]

根据 Lov\'asz 条件推广知：
\[ \|\tilde{\bm{b}}_i\|^2\leqslant2^{j-i}\|\tilde{\bm{b}}_j\|^2 \]

因此：
\[ \|\bm{b}_i\|^2\leqslant2^{i-1}\cdot2^{j-i}\|\tilde{\bm{b}}_j\|^2=2^{j-1}\|\tilde{\bm{b}}_j\|^2 \]

开方可得：
\[ \|\bm{b}_i\|\leqslant2^{(j-1)/2}\|\tilde{\bm{b}}_j\| \]



\subsection{对任意 $1\leqslant i\leqslant n,\ \lambda_i(\Lambda)\leqslant2^{(i-1)/2}\|\tilde{\bm{b}}_i\|$}
前 $i$ 个基向量 $\bm{b}_1,\ \bm{b}_2,\ \ldots,\ \bm{b}_i$ 是格中 $i$ 个线性无关的向量，它们的最大长度为：
\[ \max\limits_{1\leqslant k\leqslant i}\|\bm{b}_k\| \]

由于对任意 $1\leqslant i\leqslant j\leqslant n$ 有：
\[ \|\bm{b}_i\|\leqslant2^{(j-1)/2}\|\tilde{\bm{b}}_j\| \]

所以对任意 $k\leqslant i$ 有：
\[ \|\bm{b}_k\|\leqslant2^{(i-1)/2}\|\tilde{\bm{b}}_i\| \]

由于第 $i$ 个连续极小值 $\lambda_i(\Lambda)$ 的定义是被 $i$ 个线性无关向量的最大长度所上界的，因此：
\[ \lambda_i(\Lambda)\leqslant\max\limits_{1\leqslant k\leqslant i}\|\bm{b}_k\|\leqslant2^{(i-1)/2}\|\tilde{\bm{b}}_i\| \]



\subsection{对任意 $1\leqslant i\leqslant n,\ \lambda_i(\Lambda)\geqslant2^{-(n-1)/2}\|\bm{b}_i\|$}
设 $\bm{v}_1,\ \bm{v}_2,\ \ldots,\ \bm{v}_i$ 是取得连续极小值 $\lambda_1,\ \lambda_2,\ \ldots,\ \lambda_i$ 的 $i$ 个线性无关向量，将它们在正交基下展开，由于它们线性无关，必然存在至少一个向量 $\bm{v}^{*}$，其展开式的最大非零项索引 $j\geqslant i$，因此：
\[ \lambda_i(\Lambda)=\max\|\bm{v}_k\|\geqslant\|\bm{v}^{*}\|\geqslant\|\tilde{\bm{b}}_j\| \]

由于对任意 $1\leqslant i\leqslant j\leqslant n$ 有：
\[ \|\bm{b}_i\|\leqslant2^{(j-1)/2}\|\tilde{\bm{b}}_j\| \]

所以：
\[ \|\bm{b}_i\|\leqslant2^{(j-1)/2}\|\tilde{\bm{b}}_j\|\leqslant2^{(n-1)/2}\|\tilde{\bm{b}}_j\| \]

因此：
\[ \lambda_i(\Lambda)\geqslant\|\tilde{\bm{b}}_j\|\geqslant2^{-(n-1)/2}\|\bm{b}_i\| \]



\section{证明 $B\leqslant c_1(d)N^{\tfrac{2}{d(d+1)}}$ 时导致 $\left|a_iB^i\right|<\dfrac{N}{d+1}$}
在 Coppersmith 方法中，对次数为 $d$ 的首一多项式 $f(x)$，要求解 $f(x_0)\equiv0$ 在模 $N$ 意义下有小根 $|x_0|\leqslant B$，构造如下 $d+1$ 维格 $\mathcal{L}$，其基矩阵为下三角矩阵，对角元为：
\[ \bigl(N,\ NB,\ NB^2,\ \ldots,\ NB^{d-1},\ B^d\bigr) \]

\paragraph*{第一步：计算格的行列式}~{}
\[ \det(\mathcal{L}) = N^d \cdot B^{0+1+2+\cdots+(d-1)+d} = N^d \cdot B^{\tfrac{d(d+1)}{2}} \]



\paragraph*{第二步：应用 LLL 短向量界}~{}

设 $\bm{v} = (a_0, a_1 B, a_2 B^2, \ldots, a_d B^d)$ 为 LLL 算法输出的第一个基向量，则：
\[ \|\bm{v}\|_2 \leqslant 2^{d/4} \cdot \det(\mathcal{L})^{\tfrac{1}{d+1}} = 2^{d/4} \cdot N^{\tfrac{d}{d+1}} \cdot B^{\tfrac{d}{2}} \]



\paragraph*{第三步：代入 $B$ 的界}~{}

设 $c_1(d) = \dfrac{1}{(d+1)^{2/d} \cdot 2^{1/2}}$，当 $B \leqslant c_1(d)\, N^{\tfrac{2}{d(d+1)}}$ 时，有：
\[ B^{\tfrac{d}{2}} \leqslant c_1(d)^{\tfrac{d}{2}} \cdot N^{\tfrac{1}{d+1}}= \dfrac{1}{(d+1)\cdot 2^{d/4}} \cdot N^{\tfrac{1}{d+1}} \]

代入上界：
\[ \|\bm{v}\|_2 \leqslant 2^{d/4} \cdot N^{\tfrac{d}{d+1}} \cdot \dfrac{N^{\tfrac{1}{d+1}}}{(d+1)\cdot 2^{d/4}} = \dfrac{N}{d+1} \]

当 $B < c_1(d)\, N^{\tfrac{2}{d(d+1)}}$严格不等号时，上式为严格小于，其中：
\[ c_1(d)=\dfrac{1}{\sqrt{2}(d+1)^{\tfrac{2}{d}}} \]



\paragraph*{第四步：逐分量的界}~{}

由于 $\left|a_iB^i\right| \leqslant \|\bm{v}\|_\infty \leqslant \|\bm{v}\|_2$，故对所有 $i = 0, 1, \ldots, d$ 有：
\[ \left|a_iB^i\right|<\dfrac{N}{d+1} \]



\section{证明 Babai 算法可以解出离 $\bm{t}$ 最近的格点}
令 $\bm{B}=(\bm{b}_1,\ \bm{b}_2,\ \ldots,\ \bm{b}_n)$ 是格 $\mathcal{L}$ 的一组基，令向量 $\bm{t}\in\mathbb{R}^n$ 符合：
\[ \text{dist}(\mathcal{L},\bm{t})<\dfrac{1}{2}\min\limits_{i\in[n]}\|\tilde{\bm{b}}_i\|_2 \]

设 $\bm{v}^* \in \mathcal{L}$ 为最近格点，令误差向量 $\bm{e} = \bm{t} - \bm{v}^*$，则：
\[ \|\bm{e}\|_2 = \text{dist}(\mathcal{L},\bm{t})<\dfrac{1}{2}\min\limits_{i}\|\tilde{\bm{b}}_i\|_2 \]

将 $\bm{e}$ 在 Gram-Schmidt 基下展开：
\[ \bm{e}=\sum_{i=1}^n \epsilon_i \tilde{\bm{b}}_i \]

由于 $\{\tilde{\bm{b}}_i\}$ 两两正交，则：
\[ \|\bm{e}\|_2^2 = \sum_{i=1}^n \epsilon_i^2 \|\tilde{\bm{b}}_i\|_2^2 \]

对每个 $i$，有：
\[ \epsilon_i^2 \|\tilde{\bm{b}}_i\|_2^2\leqslant\|\bm{e}\|_2^2< \frac{1}{4}\min_j\|\tilde{\bm{b}}_j\|_2^2\leqslant\frac{1}{4}\|\tilde{\bm{b}}_i\|_2^2 \]

因此 $|\epsilon_i| < \dfrac{1}{2}$ 对所有 $i$ 成立。



\paragraph*{Babai 算法对 $i$ 从 $n$ 到 $1$ 归纳}~{}

设 $\bm{v}^* = \displaystyle\sum_{i=1}^n c_i \bm{b}_i$ 且 $c_i \in \mathbb{Z}$，由 Gram-Schmidt 的定义，$\bm{b}_j = \tilde{\bm{b}}_j + \displaystyle\sum_{k<j}\mu_{j,k}\tilde{\bm{b}}_k$，故 $\bm{v}^*$ 在 $\tilde{\bm{b}}_n$ 方向上的系数恰为 $c_n$。

当 $i=n$ 步时，$\bm{t}$ 在 $\tilde{\bm{b}}_n$ 方向的系数为 $c_n + \epsilon_n$，由 $|\epsilon_n| < \dfrac{1}{2}$ 得：
\[ \left\lfloor c_n+\epsilon_n\right\rceil=\mathrm{round}(c_n + \epsilon_n) = c_n  \]

算法减去 $c_n \bm{b}_n$，得到残差：
\[ \bm{t}^{(n-1)} = \bm{t} - c_n\bm{b}_n = \sum_{i=1}^{n-1}c_i\bm{b}_i + \bm{e} \]



\paragraph*{归纳步骤}~{}

设在第 $j+1,\ j+2,\ \ldots,\ n$ 步均已正确舍入，在第 $j$ 步，当前目标向量为：
\[ \sum_{i=1}^{j}c_i\bm{b}_i + \bm{e} \]

其在 $\tilde{\bm{b}}_j$ 方向的系数为 $c_j + \epsilon_j$，仍满足 $|\epsilon_j| < \dfrac{1}{2}$，故：
\[ \left\lfloor c_j + \epsilon_j \right\rceil=\mathrm{round}(c_j + \epsilon_j) = c_j \]

算法在每步均正确恢复 $c_i$，最终输出 $\displaystyle\sum_{i=1}^n c_i\bm{b}_i = \bm{v}^* $，即 $\bm{t}$ 的最近格点。


\end{document}